\section{Introduction}
A self-avoiding walk of prescribed length and a self-avoiding walk whose length $L$ is penalized by $e^{-L/N}$ need not have the same distribution in space. The first law fixes one coefficient of a generating function; the second sums all its coefficients, with $N$ setting the discount scale. A third law is obtained by joining independent irreducible bridges forever. Understanding the relations among these laws requires control of their normalizing masses, as well as estimates for the individual paths.

We study these changes of law on the regular honeycomb lattice. Realize its vertices as the centers of the triangles in an equilateral triangular tiling with side length one, and put
\[
 \rho=(2+\sqrt2)^{-1/2}.
\]
A port is the midpoint of a tiling edge on a domain boundary. A strict bridge starts at a prescribed port on one horizontal tiling line, ends at a port on a higher parallel line, and has all its visited vertices strictly between the two lines. A port path with $L$ visited vertices has critical weight $\rho^L$. The ports carry no weight, so adjoining bridges concatenate with exactly the product of their weights. The terminal height and horizontal position are free unless specified.

An intermediate horizontal line crossed exactly once determines a unique cut of a bridge. Cutting at all such lines decomposes the bridge into irreducible bridges. Their critical weights sum to one and therefore define a probability law $p$. Independent $p$-distributed pieces can be concatenated without further avoidance conditions: successive pieces occupy disjoint open slabs. This construction provides the common probability space used throughout the paper.

For an ordinary finite walk $\gamma=(\gamma_0,\ldots,\gamma_n)$, write
\[
 A(\gamma)=|\gamma_n-\gamma_0|,
 \qquad D(\gamma)=\diam\{\gamma_0,\ldots,\gamma_n\}.
\]
Uniform full-plane walks start at one prescribed vertex and have $n$ edges. Uniform half-plane walks start at the inward vertex of a prescribed boundary port and have a prescribed number of visited vertices in the strict half-plane, with their last vertex unrestricted. Thermal walks have variable length and weight $\rho^L e^{-L/N}$, normalized over all such walks. Using edge length instead of vertex length in this last law multiplies all weights by one common factor and gives the same probability measure.

Here are the principal spatial conclusions, with the law included in each statement. The notation $X_t=t^{a+o(1)}$ in probability means that for every $\varepsilon>0$ the event $t^{a-\varepsilon}\le X_t\le t^{a+\varepsilon}$ has probability tending to one.
\begin{itemize}
\item Under the infinite irreducible-bridge law, the distance of the $n$th vertex from the initial port is $n^{3/4+o(1)}$ almost surely, and the number of visited vertices within distance $R$ is $R^{4/3+o(1)}$ almost surely (Theorem~\ref{R:renew:infinite}). Both uniform and fugacity-weighted finite half-plane laws converge on each fixed initial path to this law (Theorem~\ref{R:renew:uniform} and Lemma~\ref{R:renew:halfplane}).
\item Uniform strict bridges with $n$ visited vertices have terminal distance and diameter $n^{3/4+o(1)}$ in probability at every sufficiently large even $n$. The parity restriction is part of the port convention. Critical bridges conditioned instead on height $h$ have length $h^{4/3+o(1)}$ in probability (Theorems~\ref{R:exact:bridge} and~\ref{R:exact:height}).
\item Uniform endpoint-free half-plane and full-plane walks have endpoint distance and diameter $n^{3/4+o(1)}$ in probability along a common set of lengths of natural density one. Their upper spatial bound holds at every sufficiently large length, with arbitrarily high polynomial accuracy (Theorems~\ref{R:cal:halfplane-laws} and~\ref{R:cal:plane-laws}).
\item For thermal half-plane and full-plane walks, $L=N^{1+o(1)}$ and $A,D=N^{3/4+o(1)}$ in probability. Their positive moments have the corresponding exponents: for each fixed $p>0$, $\E L^p=N^{p+o(1)}$ and $\E A^p,\E D^p=N^{3p/4+o(1)}$ (Theorems~\ref{R:thermal:law} and~\ref{R:cal:halfplane-laws}).
\end{itemize}
The full-plane endpoint lower bound at fixed length retains the density-one qualification. The thermal statement concerns every large discount scale $N$. These conclusions use different normalization arguments, described below. Nienhuis's analysis of the two-dimensional $O(n)$ model predicts the spatial exponent $\nu=3/4$ for self-avoiding walk in the $n\to0$ polymer limit~\cite[Eqs.~(1)--(2)]{Nienhuis1982}; the results here retain the sampling-law and length qualifications stated above.

\subsection*{Bridge theory and the additional estimates}
Kesten's normalization for irreducible bridges underlies the infinite-bridge construction of Madras and Slade~\cite{Kesten1963,MadrasSlade1993}. Lawler, Schramm and Werner identify this measure as the fixed-length half-space limit and describe its fugacity counterpart~\cite[Appendix]{LSW2004}. Dyhr, Gilbert, Kennedy, Lawler and Passon give explicit finite-prefix descriptions of both limits and identify the strip law by conditioning on a renewal height~\cite[Propositions~2.2--2.3, 2.9 and~2.10]{Dyhr2011}. These works use vertex conventions on the integer lattice; the strip construction of Dyhr et al. also includes a final bond from the conditioning height to the target height. We adapt the renewal arguments to the present honeycomb port convention and prove the required normalizations here. The infinite law, a fixed-height bridge law and a law conditioned on one exact length are different uses of the construction.

Duminil-Copin and Smirnov proved the value of $\rho$ using a parafermionic observable on the honeycomb lattice~\cite{DCS2012}. Beaton et al. use a honeycomb mid-edge renewal construction and prove that the critical bridge mass tends to zero as the strip height grows~\cite[Theorem~10 and Appendix]{BeatonEtAl2014}. Krachun and Panagiotis obtain a positive-power decay bound for this mass and show that a uniform $n$-step honeycomb walk remains within distance $O(n/\log n)$ of its root with probability tending to one~\cite{KrachunPanagiotis2026}.

The sharper finite estimates used below require additional strip and marked-polygon arguments. We state the companion results as estimates of positive walk weights. In Section~\ref{R:sec:renewal}, the finite bridge estimates of Theorem~\ref{R:renew:finite-input} and the renewal identities yield joint bounds for the height, length and diameter of one irreducible. Section~\ref{R:geo:section} uses further finite strip and polygon estimates to control avoidance and absolute path masses. The analytic finite estimates are proved in the cited companions; the present paper proves the renewal and geometric transfers needed when the sampling law changes.

The key exponents for one irreducible are
\[
 \alpha=\frac9{16},\qquad \beta=\frac34.
\]
Write $H,L,D$ for the height, vertex length and diameter of one irreducible. On the main route the finite estimates give
\[
 1-\E_p e^{-uL}\asymp u^\alpha,\qquad
 1-\E_p e^{-sH}\asymp s^\beta,\qquad
 \Pp_p\{D>r\}\le Cr^{-\beta}
\]
as $u,s\downarrow0$ and $r\to\infty$. Here $\asymp$ means upper and lower bounds with positive constants independent of the parameter. The ratio $\alpha/\beta=3/4$ explains the spatial scale after many pieces have been joined. Turning this observation into a theorem at one exact length requires more than a sum estimate.

\subsection*{How the normalizations are controlled}
Our main exact-length route uses these bounded-factor finite estimates. If $S_L(k)=L_1+\cdots+L_k$ for independent irreducibles, the local denominator is
\[
 u_n=\Pp_p\{S_L(k)=n\text{ for some }k\}\ge c n^{\alpha-1}
 \quad\text{for all sufficiently large even }n.
\]
Section~\ref{R:sec:exact} proves this by exponential tilting and Fourier inversion on the compatible length lattice. Fourier concentration for the remaining pieces then pays for sequences with atypical total height or diameter. The identification of the uniform half-plane prefix limit in Section~\ref{R:sec:renewal} is a separate ratio argument.

A second, self-contained route begins with a logarithmic window of bridge lengths. Its length information is weaker:
\[
 c u^\alpha(\log(1/u))^{-P}\le 1-\E_p e^{-uL}\le C u^\alpha.
\]
A Fourier difference estimate first gives a long interval on which the renewal mass has the required lower bound. Fresh independent increments then move the remaining length into that interval. The logarithmic window permits enough progress at each stage that the total landing cost is only $n^{-o(1)}$, giving $u_n\ge n^{\alpha-1-o(1)}$. Section~\ref{R:cal:section} proves this alternative under its own hypotheses, with a fixed logarithmic exponent below one and a stopping interval whose width remains large.

Section~\ref{R:geo:section} next establishes the geometric estimates used in the thermal transfer of Section~\ref{R:sec:thermal}. A half-plane walk consists of a list of complete irreducibles followed by a remainder with no further renewal. In a prescribed-prefix probability, the total weight of that remainder cancels exactly. For an unrestricted plane walk we obtain a denominator by putting a nonempty bridge between two exterior remainders. Recovering the first and last single-crossing lines makes this construction injective. To bound small endpoint displacement, we cut a path at its lowest vertex and use the two-end estimate of Corollary~\ref{R:geo:two-ends}, retaining one independent height-smoothing variable. The two remainder weights then appear in both numerator and denominator and cancel.

The censored-arm estimates in Section~\ref{R:cal:geometry-section} give a separate geometric route from the logarithmic finite inputs to absolute spatial bounds. Section~\ref{R:cal:uniform-section} then treats moments and uniform walks, where there is no comparable exact cancellation at each length. In the half-plane, disjoint long-prefix events compare masses at nearby lengths. In the full plane, we insert a bridge at a minimum and use censored survival to control the inverse multiplicity. Both arguments average logarithms of the nearby-length masses: the interior terms telescope and leave a boundary error. This is why the resulting lower laws hold on a density-one set.

Next, let $G_t(b)$ be the total weight $(\rho e^{-t})^{\ell(\gamma)}$ of ordinary walks from a fixed vertex $o$ to $b$, with edge length $\ell$. Section~\ref{R:resp:section} proves that the convergence threshold in $s$ of
$\sum_bG_t(b)e^{s e\cdot(b-o)}$ is $t^{3/4+o(1)}$, uniformly for unit directions $e$. The fixed-length free energy
\[
 f_e(s)=\lim_{j\to\infty}\frac1j\log\sum_{\ell(\gamma)=j}\rho^j e^{s e\cdot(\gamma_j-o)}
\]
exists and is $s^{4/3+o(1)}$ as $s\downarrow0$, uniformly in $e$. Its two one-sided derivatives are $s^{1/3+o(1)}$, in agreement with Pincus's force-extension scaling relation $1/\nu-1$ at $\nu=3/4$~\cite[Eq.~(I.4)]{Pincus1976}. Here the response follows from the threshold bounds and convexity; the distinction from fixed-force ballistic results is discussed in Section~\ref{R:resp:section}.

Prescribing a terminal vertex requires a further estimate: if $|b_N-o|=N^{3/4+o(1)}$, the thermal law conditioned to end at $b_N$ has length $N^{1+o(1)}$ and diameter $N^{3/4+o(1)}$ in probability, with the corresponding powers for all positive moments (Theorem~\ref{R:resp:prescribed}). Proposition~\ref{R:resp:point-local} obtains a uniform two-dimensional lattice lower bound from subsequential local limits. A lower partition bound supported in a corridor then permits this conditioning. The separate spatial-rate statement takes distance to infinity at fixed discount and then takes the discount to zero; it does not identify these two orders of limits.

The final three sections develop consequences of other finite inputs. Section~\ref{R:alt:section} retains a length reward in a logarithmically enlarged box and a completed-prefix expectation in a fixed-proportion ball. Section~\ref{R:alt:j} derives the irreducible length tail from signed cylinder and marked-polygon estimates. Section~\ref{R:alt:g:section} gives a geometric route to the infinite and thermal laws, using future-renewal averaging for the infinite path and recoverable tall-bridge insertion for the thermal path. These routes state their own finite inputs and retain the precision those inputs supply.
